\section{Limitations and Future Work}
\label{sec:limitations}

\paragraph{Contextual transferability.}
The five deployments span different public and quasi-public settings, but they all took place in New York City and involved the same sanitation-oriented service concept. The specific institutions, infrastructure, and local expectations that participants drew on are therefore tied to this municipal and service context. Future work should examine whether the recurring sensemaking questions identified here also appear in other cities, cultural settings, and public-service domains, and how the resources used to answer those questions change across contexts. Our interview corpus also provides uneven coverage of multilingual encounters because some non-English speech was not transcribed reliably. This limits the accounts represented in our analysis and points to the need for stronger multilingual support in future field studies.

\paragraph{Place-level and participant-level variation.}
Our analysis foregrounds the knowledge participants drew from the deployment setting, but participants also brought their own experiences to the encounter. Occupation, familiarity with the site, prior experience with robots and other technologies, and experience in other neighborhoods or institutions could all inform how a participant interpreted the Trashbot. The present study does not systematically separate these participant-level resources from knowledge associated with the site. Future work could examine how personal experience and place-based knowledge combine when people interpret public robots, including cases in which they point toward different conclusions.

\paragraph{Wizard-of-Oz and autonomous deployment.}
The Trashbots were remotely operated throughout the deployments, while participants initially encountered each robot as apparently autonomous. This operating arrangement allowed us to use the same platform design and operating protocol across sites, but the study captures sensemaking around perceived autonomy rather than sustained interaction with fully autonomous systems. Autonomous operation may introduce navigation failures, delays, recovery behavior, and other cues that affect how people read the robot's actions and judge its competence or appropriate role. Longer deployments may also change these interpretations as people become familiar with the robot and its institutional context. Future work should examine whether the sensemaking strategies and factors identified here persist under autonomous and longer-term deployment.